
\documentclass{article}
\usepackage{import}
\usepackage{epstopdf}
\usepackage{bsymb}
\usepackage{alltt}
\usepackage{amsmath}
\usepackage{amssymb}
\usepackage{multicol}
\usepackage{color}
\usepackage{hyperref}

\begin{document}

\title{Tasking Event-B: A Code Generation Extension for Event-B}


\author{Andrew Edmunds\\
School of Electronics and Computer Science,\\
University of Southampton, UK\\
\url{ae2@ecs.soton.ac.uk}\\
\and
Michael Butler\\
School of Electronics and Computer Science,\\
University of Southampton, UK\\
\url{mjb@ecs.soton.ac.uk}\\
}

\maketitle
%------------------------------------------------------------------------------
% Abstract
%
\begin{abstract}
Event-B is a formal approach to modelling systems, based on set theory, predicate logic and arithmetic. Various techniques are used during modelling, including refinement, and decomposition. This article describes an extension to the Event-B approach, which we call Tasking Event-B, that facilitates automatic generation of source code, from annotated Event-B models. We believe that automatic code generation makes a useful contribution to the Rodin tool-set, by contributing a link in a coherent tool-chain. Automatically generating code from the Event-B model can be seen as a productivity enhancement. It removes a source of errors, that of manually coding for each development. To validate the approach we have undertaken case studies and taken part in an industrial collaboration. We present a number of case-studies to illustrate our work, in this article.
\end{abstract}

\section{Introduction}\label{intro}
Event-B~\cite{ABR10} is one of a number of formal methods that may be used to model systems where a high degree of reliability is required. Event-B was inspired by its predecessor, \emph{Classical-B}~\cite{TheBBook}. It is a modelling language, used with a supporting tool platform, Rodin~\cite{abrial10rodin}; so named from the project in which it was developed~\cite{RodinTool}. Further work was undertaken in the DEPLOY~\cite{DEPLOY} project to assess the approach with a number of the industrial partners. To derive the greatest benefit from the formal modelling approach, it is desirable to use the formal modelling artefacts to generate implementations for, at least, some parts of the system being modelled. During the latter stages of the DEPLOY project, automatic generation from Event-B models began to evolve, as the platform stabilized, where we targeted multi-tasking, embedded control systems. Initially, we chose Ada~\cite{ada2005} as a basis for our approach. This was not only because of it's suitability for the application domain, but also because the programming constructs are well-considered programming abstractions. Ada maps well to Event-B modelling elements, which is easy for modelling, and simplifies the translation to code. We do not attempt to model all aspects of the system, e.g. time; but we model the evolution of the system state, and are able to specify properties relating to the state, that are important for system safety. In current work, in the Advance project~\cite{advance}, we are concerned with modelling and co-simulation of Cyber-physical systems. We can, again, make use of automatic code-generation techniques to provide implementations for use in co-simulation.    


\subsection{Some History}
Initial development of the code-generation approach began from an object-oriented perspective, with the implementation-level notation, object-oriented, concurrent-B (OC-B)~\cite{Edmunds2009}. The focus of this work was to use Event-B  to model and generate code for safe multi-tasking with object-oriented languages. We initially targeted the Java~\cite{JavaSpec} language. We found, however, that this approach gave rise to a notation where the semantic gap between Event-B and the OC-B was quite large. This was not optimal for developers of Event-B models, who then wanted to write implementation-level specifications, from which code could be generated. We also found that models quickly became large and intractable; we recognized the need for decomposition of models into smaller units. At that time, Event-B tools were at an early stage of development, and machine decomposition was still in development. 

We decided to modify the approach, to use the newly developed decomposition approach~\cite{decomp2010c,decomp2010b}, and use an implementation-notation more closely related to Event-B (than OC-B, which was an object-oriented notation). The focus, then, became generating safe concurrent implementations, rather than targeting object-oriented technology, per se. This gave rise to our approach, known as Tasking Event-B. Our interest in safe, concurrent implementations drove us towards the Ada tasking model. At the implementation-level, an Event-B model is effectively a detailed model of an implementation. We found that the Ada~\cite{ada2005} programming language provides clearly defined constructs, that map well to Event-B, see Fig.~\ref{fig:B_Ada}. Ada tasks and protected objects can be modelled by Event-B machines in a one-one mapping. We distinguish between the two entities by annotating the machines with the AutoTask and Shared keywords. We describe the mapping in more detail in Sects.~\ref{targets} and~\ref{}.
%
\begin{figure}
\centering
\includegraphics[width=0.5\textwidth]{graphics/B_Ada.png}
\caption{Mapping between Event-B and Ada}
\label{fig:B_Ada}
\end{figure}
%

In Sect.~\ref{eventb} we introduce the core Event-B features to the reader, and compare Event-B with some other formal approaches.  In Sect.~\ref{features} we provide an overview of some additional Event-B features required for code generation, and discuss some of the potential target programming language targets.

In Sect.~\ref{} we...
In Sect.~\ref{} we...

\section{Event-B Modelling}\label{eventb}
The formal methods related to the work presented here can be categorized as state-based formal methods. Alternative, but not unrelated, approaches are categorized as process-based methods. Classical-B~\cite{TheBBook,CNP,B4Free,atelierB} and its successor, Event-B are said to be state-based, since they focus on modelling the changes of state, not the behaviour of processes. In Classical-B, state updates are modelled by guarded operations, where the operation is an analogue of a procedure call in a programming language. In Event-B, state updates are modelled by guarded events, providing a more abstract view of the way a system evolves. Event-B can be used to model systems at an abstract level; and by adding more detail (using a technique called refinement) it can model the software aspects of systems too. Both methods are set theoretic modelling approaches that incorporate a notion of proof to show that important system properties are maintained. The former is primarily an approach to software systems development, the latter more widely applicable to system-modelling. In an effort to make modelling and proof easier, Event-B was developed to overcome some of the difficulties encountered when using in Classical-B. The main differences between Classical and Event-B are highlighted in~\cite{Hallerstede07}, and inspiration was also drawn from action systems~\cite{Back1990133}.

It is fair to say that Event-B is not just a formal modelling language; the name is used to describe both a notation, and a methodology. In addition to this a mature tool-platform called \emph{Rodin}, named after its development programme, complements the methodology. The main modelling components of Event-B are contexts and machines. Contexts are used to model static features using sets, constants, axioms, and theorems. Machines are used to model variable state using \emph{variables}. A third, more recent addition, is the Theory component; where a developer can augment the bundled mathematical language, and rule-base, with new (inference and re-write) rules, data types, and operators. During the modelling process, changes to the components result in automatic generation of proof obligations, which must be discharged in order to show that the development is consistent. The proof obligations generated in classical-B are often complex, the Event-B approach results in simpler proof obligations as described in~\cite{Hallerstede07}, since Event-B consists of a simplified action syntax, giving rise to simpler proof obligations. A further simplification was made by adopting an event-based approach, where each atomic event has a predicate guard and an action consisting only of assignment statements. Events correspond to operations in the B-method; operation specification was more expressive, and included constructs for specifying operation preconditions (as part of its Design by Contract approach), operation calls, return parameters, and more complex structures for branching and looping. These constructs are not features of Event-B. Due to these simplifications (and more efficient proof tools) a large number of the proof obligations may be discharged automatically, by the automatic provers. Where un-discharged proof obligations remain, the user has, at their disposal, an interactive prover. Various techniques can be applied, to discharge the proof obligations, such as adding hypotheses; or making use of the hyperlink-driven user interface, for rule and tactic application. In the early stages of development with Event-B, a developer will begin by abstracting, and modelling, the observable events occurring in a system. Event-B, as the name suggests, takes an event-based view of a system; where events occur spontaneously from the choice of enabled events. An event is said to be enabled when the guard is true, and the state updates, described in the event actions, can take place; otherwise it is disabled, and none of its updates can occur.  

\subsection{An Event-B Example}
An example of an Event-B machine can be seen in Fig.~\ref{fig:controllerSpec2}. It shows an abstract model of a pump controller, used in one of the case studies. We will use this model to describe some features of Event-B. But first we introduce the case study, which models a discrete \emph{pumpController}. The model describes a system where the controller receives a value of the fluid level, and a boolean value representing a user-request to turn the pump. Based on the inputs to the controller, a command to turn the pump on may be issued, or a warning is issued if a minimum level \emph{MIN} has been reached.    
%
%
%
\begin{figure}[h]
\centering
\begin{minipage}{0.8\textwidth}
\textcolor{blue}{MACHINE} m1 \textcolor{blue}{REFINES} m0 \textcolor{blue}{SEES} ctx \\
\textcolor{blue}{VARIABLES} \text{m\_level, c\_level, e\_level, m\_pumpOnReq, c\_pumpOnReq, e\_pumpOnReq,} \hspace*{0.2cm} m\_pumpOnCmd, c\_pumpOnCmd, e\_pumpOnCmd, m\_warn, c\_warn, e\_warn,\\
\hspace*{0.2cm} c\_level\_internal, c\_pumpOnReq\_internal\\
\textcolor{blue}{INVARIANTS}\\
\hspace*{0.2cm}(c\_level\_internal $\leq$ MIN $\land$  c\_pumpOnReq\_internal = TRUE $\limp$  c\_warn = TRUE)\\
\hspace*{0.2cm} $\land$ (c\_level\_internal $>$  MIN $\land$  c\_pumpOnReq\_internal = TRUE\\
\hspace*{0.5cm} $\limp$  c\_pumpOnCmd = TRUE)\\
\hspace*{0.2cm} $\land$ (c\_level\_internal $\in  \intg$)\\
\hspace*{0.2cm} $\land$ (c\_pumpOnReq\_internal $\in$  BOOL) \ldots\\
\textcolor{blue}{EVENTS}\\
\textcolor{blue}{INITIALISATION} c\_level :=  100 $\pprod$ m\_level := 80 $\pprod$ c\_pumpOnReq :=  FALSE $\pprod$ \ldots\\
\textcolor{blue}{EVENT} fmiSetBoolean\_c \textcolor{blue}{REFINES} fmiSetBoolean\_c\\
\hspace*{0.2cm}\textcolor{blue}{ANY} p\\
\hspace*{0.2cm}\textcolor{blue}{WHERE} p = c\_compound $\land$ p $\in$ BOOL  \\
\hspace*{0.2cm}\textcolor{blue}{THEN} m\_pumpOnCmd :=  p\\
\hspace*{0.2cm}\textcolor{blue}{END}\\
\ldots
\end{minipage}
\caption{An Event-B  Pump Controller Model}
\label{fig:controllerSpec2}
\end{figure}
%
%
%
\begin{figure}[b]
\centering
\begin{minipage}{0.3\textwidth}
\textcolor{blue}{CONTEXT} ctx  \\
\textcolor{blue}{CONSTANTS} MIN\\
\textcolor{blue}{AXIOMS} MIN = 10
\end{minipage}
\caption{An Example Context}
\label{fig:context}
\end{figure}
%
%
 In Fig.~\ref{fig:controllerSpec2}, we see that machine \emph{M1} refines another machine \emph{M0}; we will discuss refinement in Subsect.~\ref{ref}. It also has a \emph{SEES} clause; this makes the contents of a context visible to a machine. Contexts may contain sets, constants, axioms and theorems, and example context can be seen in Fig.~\ref{fig:context}. There are variables representing the internal state of the controller, and invariants providing type information for variables. Invariants are also used to describe the safety properties of the system. This describes a required safety property, that if the level is at or below \emph{MIN}, and a user's pump-on request is detected, then a warning will be issued. Also, if the level is OK and a pump-on is requested, then the state \emph{pumpOnCmd = TRUE} is  set.  Following the \emph{INVARIANTS} clause are the model's \emph{Events}. The \emph{Initialisation} event is special event, since it has no guards. The initialisation event of a machine must occur before any other event in the machine is enabled. The event in the figure has a parameter \emph{p}, in the \emph{ANY} clause. Parameters can be used to represent information flow, in and out of events, or they can represent a \emph{local} variable within the scope of the event. The event guard is defined in the \emph{WHERE} clause, in the example, where \emph{p} is typed as a Boolean. The guard relates the parameter to a machine variable \emph{c\_pumpOnCmd}, in the predicate $ p = c\_compound$. The event action appears in the \emph{THEN} clause, where the parameter is assigned to the variable \emph{m\_pumpOnCmd}, in the expression $m\_pumpOnCmd := p$.


\subsection{Refinement and Extension}\label{ref}
As we mentioned earlier in the section, Event-B makes use of a technique called refinement, where a machine can be refined by another. Fig.~\ref{fig:refine} shows this relationship.
%
\begin{figure}
\centering
\includegraphics[width=0.5\textwidth]{graphics/refinement.png}
\caption{Refinement and Extension}
\label{fig:refine}
\end{figure}
%
%
The refined machine is augmented with state variables, events and invariants, to provide a more detailed specification satisfying the properties, specified in the invariants, of the abstract specification. The counterpart to refinement of machines, is extension of contexts. It is then possible to build upon pre-existing contexts, using the \emph{EXTENDS} clause, by adding more sets, constants, axioms and theorems. When a machine \emph{SEES} a context, the contexts that it extends are also accessible to the machine. In a refinement, new variables, events and invariant properties can be added. Existing events can be modified, but in a restricted manner. Machine refinement is transitive and leads to a hierarchical structure. Refinements are related to their more abstract counterparts in such a way that, a valid refinement always satisfies the specifications higher in the refinement hierarchy. In this way, important system properties can be specified at a high level of abstraction, and maintained down through the refinement chain. The Event-B tools are responsible for generating the proof obligations relating to refinement; these must be discharged in a similar way to those generated for proof of machine consistency. It is often necessary to specify a linking invariant, to describe the relationship between the variables of the abstract and refinement machines. Inspection of the proof obligations can assist in this task since some of the un-discharged proof obligations provide information about this link. In some cases we may model entities in an abstraction that are defined in the event parameters; and in the refinement these entities may be introduced to the model as machine variables. To assist with the proof effort we have a slightly different strategy to that of refining abstract variables with concrete variables. We link the parameters of the abstract event, with their concrete counterpart, using a $WITNESS$. This construct is a predicate describing the relationship between an event parameter (that disappears) from an abstract model, and the corresponding (refining) variable in the concrete model. Another feature of Event-B is the ability to  refine one atomic event with a number of events, thus breaking the atomicity, as described in~\cite{Butler08}. Eventually, at the end of a refinement chain the models are detailed enough to accurately describe an implementation. But Event-B is a modelling language, and there is a disjunction between the description of the system in Event-B, and commonly used programming languages, such as Java~\cite{JavaSpec}, C~\cite{KernighanR88} and Ada~\cite{ada2005}. Addressing the semantic gap between Event-B and programming constructs is a contribution of this article. 
%


\section{More about Event-B}

\subsection{Refinement}

\subsection{Decomposition and Composition}

\subsection{Theories}

\subsection{ProB}


\section{Tasking Event-B}\label{TEB}
In this section we begin our presentation of Tasking Event-B. We begin by describing the Event-B semantics of the notation. An Event-B model is an abstraction of a system, the evolution of the state is described using events. The guard of an event describes the conditions which must hold prior to an event being enabled. The actions describe the values of the state variables after the event has occurred. Proof obligations are generated, to ensure that the actions update the state with values that satisfy invariant. As a result, a system that is consistent, can be described using events that do not impose an ordering: since in Event-B any enabled event can occur.

Now, in an implementation-level specification wish to reason about task-like behaviour, with implementation constructs, such as sequencing. So we introduce a notation for the purpose of ordering of events. The notation can be used to assist with code generation, since it allows the developer to introduce an ordering where none might exist. In the case that such an ordering does exist in the abstraction, the developer can generate a new Event-B model from the notation, and show that it a refinement of the more abstract one.  When considering translation to code, it is usually necessary to work with a subset of implementable Event-B constructs. We consider \emph{implementable constructs} to be those that are available in (or map well to) a programming language. We would therefore usually not consider non-deterministic assignment to be implementable, for instance, and add a restriction; that these are `refined out' of the implementation-level model. Annotations are added to both machines and contexts. As well as assisting with code generation, the annotations are used to generate an Event-B model of the implementation. In Subsect.~\ref{flow} we introduce the notation for specifying control flow, and provide Event-B, and implementation, semantics for this. In Subsect.~\ref{tasks} we discuss how we relate tasks, and protected objects, to Event-B machines.

In the following discussion, we distinguish between a task's behaviour, such as it's life-cycle (which is not modelled formally) necessary for implementation; and the updates to state (which are modelled formally). Our notation allows specification of both, formal and non-formal, aspects of a specification. An example of a non-formal aspect, is the task type which might be periodic or one-shot. Specifying that a task should repeat every so often, or happen just once. 



 
\subsection{Flow Control for Events}\label{flow}
To enable us to impose an order on events we introduce \emph{Tasking Event-B}, a textual extension to standard Event-B. The first notion we will look at is sequential ordering, for which we introduce a semi-colon `;' sequence operator. In a model with two events \emph{evt1} and \emph{evt2}, we can write \emph{evt1;evt2} to specify a sequential ordering. We can provide Event-B semantics for the ordering, by introducing an abstract program counter to the model. An enumeration of constants models the abstract program counter values; one per event. An abstract program counter variable models the currently enabled event. Guards enforce the ordering, and actions update the abstract program counter. An example can be seen in Fig.~\ref{fig:seq}, where the program counter constants are \emph{evt1} and \emph{evt2}, \emph{pc} has type \emph{pcValues}, and the program counter enumeration is defined by a partition, using the Event-B partition operator, so $partition(pcValues, \{evt1\}, \{evt2\},\{term\})$. The partition means that the values of the enumeration are distinct, i.e. $evt1 \neq evt2$ and so on. \emph{term} is label indicating a final state where no event is enabled.
%
\begin{figure}
\centering
\begin{minipage}{0.7\textwidth}
Variables evt1, evt2\\
Invariant pc $\in pcValue$ \\
Initialisation = pc $\bcmeq$ evt1 \\
evt1 = \textbf{WHEN}~ pc = evt1~ \textbf{THEN}~$A_1~ \pprod$ pc $\bcmeq$ evt2 \textbf{END}\\
evt2 = \textbf{WHEN}~ pc = evt2~ \textbf{THEN}~$A_2~ \pprod$ pc $\bcmeq$ term \textbf{END}\\
\end{minipage}
\caption{Sequence}
\label{fig:seq}
\end{figure}
Initially $evt1$ is enabled and $evt2$ is disabled, since $pc = evt1$. The update action $A_1$ occurs, with the program counter being set \emph{evt2}. This, in turn, enables $evt2$, and so on.
\subsection{Translating Sequences of Events}
We will now define some translation functions for translating Tasking Event-B to Event-B. These functions effectively add program counter guards, and actions to the event. In the first instance, the abstract syntax of TEB ::= EVENT $|$ SEQ, where the concrete syntax of a sequence \emph{SEQ} is a semi-colon `;'. When translating sequences of events, we assume that events have no guards other than the program counter guards. To add the guards and actions modelling a program step, seen in Fig.~\ref{fig:seq}, we define a translation $TEB\_{pc}$. It takes, as parameters, a Tasking Event-B construct, and the next program counter name. It returns a set of events. The $TEB_{pc}$ function is polymorphic on its inputs, with a different function application for each Tasking Event-B construct.
%
\begin{equation}
TEB_{pc} \in  TEB \cprod Name \rightarrow POW(EVENT)
\notag
\end{equation}
%
We use a function \emph{pcName} to extract an event name from an event; the name is used as a program counter variable. When $TEB_{pc}$ is applied to a sequence, it returns two events, with $TEB_{pc}$ applied to the individual events. 
%
\begin{equation}
\begin{split}
&TEB_{pc}(e_1 ; e_2, x) \leadsto \\
\{~&TEB_{pc}(e_1, pcName(e_2)),\\
&TEB_{pc}(e_2, x)~\}
\end{split}
\notag
\end{equation}
%
The $TEB_{pc}$ function is applied to each event in the returned set. We now represent an event, with guards and actions, using the notation $g \rightarrow a$; where $g$ is the guard, and $a$ is the action. 
\begin{equation}
\begin{split}
&TEB_{pc}(e_1, x) \leadsto \\
\{~&e_1 \triangleq pc = pcName(e_1) \rightarrow ( a_1 \pprod pc \bcmeq x )~\}
\end{split}
\notag
\end{equation}
This returns an updated event, with the additional guards and actions modelling a program pc.

The translation to code for a sequence is relatively straightforward. We simply map the concrete sequence operator to a statement delimiter, expand the event actions, and add a terminating delimiter in the appropriate place. One further consideration is the translation of parallel assignments to sequential, where care must be taken to translate variables on the right hand side of assignments to a local variable representing the initial values of those variables. First we define a translation function that maps from Tasking Event-B to program statements in our common language abstraction.   
\begin{equation}
TEB_{clm} \in  TEB \rightarrow ProgramStatements
\notag
\end{equation}
To resolve the mapping from parallel to sequential we define a function which generates program statements from the event $ACTIONS$. 
\begin{equation}
TEB_{iniValSubs} \in  ACTIONS \rightarrow ProgramStatements
\notag
\end{equation}
In the $TEB_{iniValSubs}$, function we translate assignments of the form $l \bcmeq E(V)$ to program assignments. Here, $l$ is a single variable identifier on the lhs of the assignment, and  E(V) represents expressions involving a set of variables $V$ on the right hand side. For each variable $v \in V$ we insert in the translated statements, a local variable of the same type as v.It is initialised so that $v_{ini} \bcmeq v$. We replace occurrences of $v$ in E(V) with $v_{ini}$. The translation function applied to the an event is therefore defined as a further translation,
%
\begin{equation}
\begin{split}
&TEB_{clm}(e_1) \leadsto \\
&TEB_{iniValSubs}(a_1)
\end{split}
\notag
\end{equation}
%
As an example, $a_1$ may be a parallel assignment, where $x \bcmeq y \pprod y \bcmeq x$. To translate this we have,
%
\begin{equation}
\begin{split}
&TEB_{iniValSubs}(x \bcmeq y \pprod y \bcmeq x) \leadsto\\
&y_{ini} \bcmeq y;~ x_{ini} \bcmeq x;~ x \bcmeq y_{ini};~ y \bcmeq x_{ini}  
\end{split}
\notag
\end{equation}
%  
\subsection{Branching}
To introduce branching to Tasking Event-B we consider the simple case first, where we have two events $e_1 = g \rightarrow a_1$ and $e_2 = \lnot g \rightarrow a_2$. We add to the syntax, IF $e_1$ ELSE $e_2$ END. The translation function takes the branch as an argument, adds the program counter information, and returns the set of updated Events. We assume that a proof obligation is generated, whereby we can show that the guards are disjoint and cover all cases.
%
\begin{equation}
\begin{split}
&TEB_{pc}(IF~ e_1~ ELSE~ e_2,~  x) \leadsto\\
& \{~ e_1 \triangleq (g_1 \land pc = pcName(e_1)) \rightarrow (a_1 \pprod pc \bcmeq x),\\
&\hspace*{0.3cm}e_2 \triangleq (\lnot g_1 \land pc = pcName(e_1)) \rightarrow (a_2 \pprod pc \bcmeq x) \}
\end{split}
\notag
\end{equation}
%
In the branching case, both events share the same enabling program counter. Instead the ordering is determined by the guard $g_1$ or $\lnot g_1$, and both branches take the same next program counter value. 

The translation to the common language model, is as follows.
\begin{equation}
\begin{split}
&TEB_{clm}(IF~ e_1~ ELSE~ e_2) \leadsto\\
& if(g_1)\{ TEB_{iniValSubs}(a_1) \} \\
& else\{TEB_{iniValSubs}(a_2) \}
\end{split}
\notag
\end{equation}

\subsection{Looping}
The Tasking Event-B while loop is written \emph{WHILE e1 END}. This simple loop repeats \emph{e1} while the guard is true. It has the following Event-B semantics for which we provide the translation,

\begin{equation}
\begin{split}
&TEB_{pc}(WHILE~ e_1,~  x) \leadsto\\
& \{~ e_1 \triangleq (g_1 \land pc = pcName(e_1)) \rightarrow (a_1 ),\\
&\hspace*{0.3cm}e_2 \triangleq (\lnot g_1 \land pc = pcName(e_1)) \rightarrow (pc \bcmeq x) \}
\end{split}
\notag
\end{equation}

The translation to the common language model, is as follows.
\begin{equation}
\begin{split}
&TEB_{clm}(WHILE~ e_1 ) \leadsto\\
& while(g_1)\{ TEB_{iniValSubs}(a_1) \}
\end{split}
\notag
\end{equation}

We will not explicitly give details of the translation of Tasking Event-B constructs from the CLM to the implementations in target languages, since we anticipate that the CLM constructs represent implementation constructs in an `obvious' way. 

\subsection{Tasks and Shared Machines}\label{tasks}
In our implementations we wish to describe two kinds of behaviour. The first is the behaviour of the task-like scheduling of activities (described by the ordering of events) specified in the Task Body, using the Tasking Event-B constructs introduced earlier in the section. The second is the behaviour of shared machines, which provide a means to share information between the tasks. Shared machines have no flow control specification of their own (i.e. no event ordering). They instead, rely on the synchronization of events arising from the decomposition, which is recorded in the composed machine structure. The events in the shared machine synchronize with the corresponding events in tasks, and contains enough information to construct an implementation. 

Tasking machines therefore consist of an Event-B machine with a task-body annotation describing the event ordering. There are two types of tasking machines, one describing the controller tasks in the system; the other describing the environment. These are known as \emph{Autotask} and \emph{Environ} machines respectively. The main difference between a model of a controller, and a model of the environment, is that it is intended that a controller task is refined down to an implementation which can be deployed on the target system. The environment model is usually an abstract model of the environment. The generated code can be discarded if required, although it is possible to generate a Java-style interface, or re-use the generated Ada procedures. The interface (or modified Ada procedures) can be implemented with calls to software driver APIs allowing the controller to interact with the real environment. The following table relates the Tasking Event-B construct to its implementation.
%
\begin{center}
\begin{tabular}{l|l}
Tasking Event-B & Implementation \\ 
\hline
AutoTask Machine & Deployable Task\\
Shared Machine &	Protected Object\\
Environ Machine &	Simulation Task and Interface
\end{tabular}
\end{center}
%
In the CLM we have a `generic' autoTask element, derived from the Tasking Event-B aouttask. The actual implementation may be in any language, but we define the semantics in terms of the Ada task, when specified in an Ada main procedure body. Ada tasks defined in this way start as soon as possible, near the beginning of a program execution. But no explicit ordering is imposed on their start-up, and  the number of tasks is fixed, determined by the number of machines in the development. The autotask might also be implemented as a Java thread, or a POSIX pthread in C; in this case the implementation should reflect that of the Ada implementation as far as possible.

 The autoTask construct, in the CLM, has a task-body containing an expansion of the tasking constructs, obtained from the application the previously introduced translation rules. The task body may contain sequences, and branches, etc. but will still contain the original Event-B predicates and expressions. Translation of these will take place in the second step, which involves the use of the theory plug-in. Ultimately autotasks are used to model a system's controllers, and translations from them will be deployed on the target system.

In the same way as autotasks are related to Ada tasks, Shared Machines are related to Ada protected objects. The protected objects provide mutually exclusive access to private data. In our translation all machine variables are translated to private variables, so it is necessary to use the object's procedures to access the data. Protected objects are owned by a task, and used to share information between tasks. Implementations in other languages may need to provide thier own mutual exclusion mechanisms, such as using Java's synchronized constructs. 

As we have already stated, Environ Machines model the environment, and their implementation is similar to an autotask's, with the exception that we allow direct communication between autotasks, and environ tasks. This is possible because we implement Ada tasks for the purposes of simulation; the environment tasks can have Ada entries. These are used in the rendezvous communication style. This allows controller implementations to poll, and set, environment values using an entry call. Ada is unusual in that it has this mechanism as a language construct. Similar implementations should be achievable in other languages.

\subsection{Theories for TEB}

\subsection{State-machines}

\section{Tooling}

\subsection{The Rodin Platform and Eclipse}

\subsection{IL1/CLM}

\subsection{Templates}

\subsection{Interfaces}

\section{Conclusions}
\subsection{Discussion}

\subsection{Related Work}
As we mentioned in the introduction, the main driver for this work has been to derive the greatest benefit from the formal modelling approach, by making use of the formal modelling artefacts to generate implementations. There are many formal notations and many have support for automatic code-generation.

The Classical-B~\cite{TheBBook} approach made use of an implementation-level notation called \emph{B0}, described in~\cite{B0RefMan}. B0 is similar to a programming language, and consists only of concrete programming constructs. These constructs map to programming constructs in a number of programming languages. \emph{B0} forms part of the Classical-B refinement chain, so the implementation-level specification refines the abstract development. Translators are available targeting programming languages such as C~\cite{KernighanR88}, and $High~ Integrity~ Ada$ (based on $SPARKAda$~\cite{SPARKAda}). 

The \emph{Z-notation} is a state-based specification notation (actually a distant ancestor of Event-B).  

The VDM-SL is a state-based formal specification language, related to $Z$, and has tool support for automatic code generation. 

VDM++ is an object-oriented extension to Models can be described textually; or using a graphical interface using UML diagrams, in much the same way as UML-B does for B and Event-B. VDM++ can be used with the VDM++ Toolbox to generate C++ and Java code. VDM++ can be used to model and implement developments with concurrently executing processes, using threads.

\emph{Object-Z}~\cite{GSmith2000} is an object-oriented approach to development using $Z$. A route to implementation is described using a translation to \emph{PerfectDeveloper}~\cite{PD} in~\cite{Stevens06}. \emph{PerfectDeveloper}'s approach is to use verified-Design By Contract, where verification conditions are generated from a specification using constructs such as pre and postconditions, class and loop invariants, and assertions. The verification conditions must be shown to hold in order to show the specified contracts are satisfied by the implementation. They are generated for each method entry to show that the precondition holds, for each method exit to show that the postcondition holds, and wherever an assertion appears. \emph{PerfectDeveloper} provides automatic, and semi-automatic, translations to Java and C++ but appears not to support concurrent processing.

There are some combined formal approach where code-generation has been investigated. CSP~\cite{HO85CSP,Roscoe1997} is a process algebraic approach to system specification in which the ordering of events occurring in a system play a major role. CSP specifications have been translated into Java code using JCSP~\cite{JCSPNet,JCSPMulti}. A hybrid CSP and Classical-B approach~\cite{SchneiderT02,SchneiderT05} combines the benefits of modelling, using both process-based and state-based techniques.  In this approach, called $CSP\pprod B$, specifications are used to constrain the order that the state-changing operations may occur; and to specify points at which the processes may interleave. The B operations synchronize with CSP events with the same name, and provide an ordering of the occurrence of B operations. ProB~\cite{LeuschelB08} is an animator and model checker for Classical-B and Event-B, and can be used with the $CSP\pprod B$ combination. It is used in the JCSProB~\cite{YangPop2007} approach, which combines CSP and Classical-B, and has a code generator inspired by JCSP. Another combined approach, amenable to model-checking, is \emph{Circus}~\cite{WoodcockC01}, which is combination of CSP and \emph{Z-notation}~\cite{Spivey89}. CSP is used to order the $Z$ operations. It can be translated to Java as described in~\cite{FreitasC06} and makes use of the JCSP library.



%------------------------------------------------------------------------------
% Refs:
%
\label{sect:bib}
\bibliographystyle{plain}
\bibliography{MyBibTex}
\end{document}

% EOF
